% ===================================================================
%  圖表第 6 號 — 屋入面、傢俬、食嘢、照明。
%
%  Traduction, faite en 2026, en cantonais ecrit d'aujourd'hui, en
%  caracteres traditionnels, sur l'ido de texto/io. Regles de la
%  colonne : voir 10-toubiu-01.tex. Cinq pieces, cinq numerotations :
%  les cles portent c1 a c5.
%
%  LE PLUS LONG TABLEAU DU LIVRET, ET CELUI OU LE CANTONAIS COMPTE SES
%  CLASSIFICATEURS. Comme le vietnamien, il en met un devant presque
%  chaque objet ; mais il n'a pas les memes, et il en a un que le
%  mandarin n'emploie pas :
%
%    個  ce qui est rond ou general   個煲, 個碗, 個櫃
%    張  ce qui est plat et large     張檯, 張床, 張氈
%    把  ce qui a un manche           把刀, 把遮, 把梳
%    條  ce qui est long              條毛巾, 條褲, 條匙
%    隻  ce qui va par paire          隻杯, 隻碟, 隻鞋
%    嚿  la MASSE informe             嚿餅, 嚿肉 — inconnu de Pekin
%
%  « 嚿 » est proprement cantonais et n'a pas d'equivalent au nord :
%  il compte ce qui n'a pas de forme, et c'est lui qui sert pour le
%  pain de savon et la motte de beurre.
%
%  LA CUISINE EST L'ENDROIT OU HONG KONG ET PEKIN SE SEPARENT LE PLUS,
%  ET LE MOT LE PLUS VOYANT EST « 雪櫃 » — mais il n'est pas ici,
%  puisque la glaciere n'existe pas dans cette maison de 1926. Ce qui
%  y est se dit quand meme autrement :
%
%    煲    la marmite   (Pekin : 锅)      鑊    le wok, la poele
%    煙通  le tuyau     (Pekin : 烟囱管)  火水燈 la lampe a petrole
%    梘    le savon     (Pekin : 肥皂)    抹枱布 le torchon
%    枕頭  l'oreiller                     被    la couverture
%
%  « 梘 » est le plus ancien : c'est le mot du savon a Canton depuis
%  le XIXe siecle, et Pekin ne l'ecrit jamais.
%
%  ET LA SALLE DE BAIN A LE MEME EMBARRAS QU'EN VIETNAMIEN. La maison
%  chinoise de 1926 n'avait pas de piece pour se baigner ; « 浴室 »
%  est un compose transparent, et le chauffe-eau (2) porte encore le
%  nom qu'il avait en arrivant.
% ===================================================================

%%K t06-apar-1 apar
\VUsaut{6.0mm}
\VUtitre{15.0pt}{60}{圖表第 6 號}
\VUsaut{1.4mm}
\VUtitre{11.4pt}{0}{屋入面、傢俬、食嘢、}
\VUsaut{0.4mm}
\VUtitre{11.4pt}{0}{照明。}
\VUsaut{2.2mm}

%%K t06-apar-2 apar
間屋起好喇。約翰阿叔搬咗入新居，我哋已經知道佢點樣分佈。而家我哋睇下佢點樣擺
傢俬。我哋先去睇：

%%K t06-tit-1 sub
\VUpk{10.2pt}{40}{睡房。}
\VUsaut{3.0mm}

%%K t06-c1-01-1 p
1. --- 我哋嚟得唔啱時候，因為個房務女工執緊\VUgras{間房}。

%%K t06-c1-02-1 p
2. --- \VUgras{梳妝檯} \textsuperscript{(1)} 上面周圍擺住啲用嚟洗面、梳頭，總之
用嚟梳洗嘅嘢。就係個\VUgras{面盆} \textsuperscript{(2)} 同個\VUgras{水壺}
\textsuperscript{(3)}、把\VUgras{頭刷} \textsuperscript{(4)}、把\VUgras{梳}
\textsuperscript{(5)}、一嚿\VUgras{梘} \textsuperscript{(6)} 同一嚿\VUgras{海綿}
\textsuperscript{(7)}，最後係一個裝漱口水嘅\VUgras{樽仔} \textsuperscript{(8)}。

%%K t06-c1-03-1 p
3. --- 地下、梳妝檯前面，就係嗰個裝污水嘅\VUgras{桶} \textsuperscript{(9)}。

%%K t06-c1-04-1 p
4. --- 我哋喺\VUgras{前廳} \textsuperscript{(10)} 見到幾個\VUgras{衣鈎}
\textsuperscript{(13)} 同兩把\VUgras{遮} \textsuperscript{(11)}，插喺個\VUgras{遮筒}
\textsuperscript{(12)} 度。

%%K t06-c1-05-1 p
5. --- 門另一邊係個\VUgras{鏡櫃} \textsuperscript{(14)}，入面裝住\VUgras{被鋪}
\textsuperscript{(15)}。

%%K t06-c1-06-1 p
6. --- 趁\VUgras{房務女工} \textsuperscript{(6)} 執緊\VUgras{張床} \textsuperscript{(17)}，
我哋睇下佢有咩部分。\VUgras{床架} \textsuperscript{(20)} 上面托住一張好嘅
\VUgras{彈簧褥} \textsuperscript{(21)}，玉婷一陣就會喺上面鋪張羊毛\VUgras{褥}
\textsuperscript{(22)}。跟住佢會由啲椅上面攞兩張\VUgras{床單} \textsuperscript{(23)}
同張\VUgras{被} \textsuperscript{(24)}。之後佢會喺床頭擺條\VUgras{長枕}
\textsuperscript{(25)} 同個\VUgras{枕頭} \textsuperscript{(26)}——呢兩樣而家同張
\VUgras{羽絨被} \textsuperscript{(27)} 一齊擺喺\VUgras{床邊地氈} \textsuperscript{(32)}
上面。佢用支雞毛掃掃走\VUgras{窗簾} \textsuperscript{(19)} 同\VUgras{床帳}
\textsuperscript{(18)} 嘅塵，咁就做完一部分嘢喇。

%%K t06-c1-07-1 p
7. --- 跟住佢仲要執下張\VUgras{床頭櫃} \textsuperscript{(28)}，上下個\VUgras{鬧鐘}
\textsuperscript{(29)} 嘅鏈，準備好\VUgras{夜燈} \textsuperscript{(30)}，攞走支
\VUgras{蠟燭} \textsuperscript{(31)} 去抹燭台。

%%K t06-tit-2 sub
\VUsaut{5.0mm}
\VUpk{10.2pt}{40}{針線籃同壁爐嘅用具。}
\VUsaut{3.0mm}

%%K t06-c1-08-1 p
8. --- 個\VUgras{針線籃} \textsuperscript{(34)} 就喺張\VUgras{凳仔}
\textsuperscript{(33)} 隔籬，一樣好需要執下。佢要摺埋幅\VUgras{刺繡}
\textsuperscript{(35)}，將啲嘢收入個\VUgras{針線檯} \textsuperscript{(38)} 度——個
\VUgras{頂針}、啲\VUgras{線} \textsuperscript{(36)}、啲\VUgras{針}
\textsuperscript{(37)} 同把\VUgras{較剪} \textsuperscript{(39)}——再用條
\VUgras{鎖匙} \textsuperscript{(40)} 鎖埋個檯蓋。

%%K t06-c1-09-1 p
9. --- 房中間張\VUgras{獨腳檯} \textsuperscript{(41)} 上面，玉婷會換咗個
\VUgras{水樽} \textsuperscript{(42)} 嘅水。佢會將\VUgras{糖}放入個\VUgras{糖罌}
\textsuperscript{(43)}，抹乾淨把\VUgras{糖鉗} \textsuperscript{(44)}，仲會收起把
\VUgras{衣刷} \textsuperscript{(46)}。

%%K t06-c1-10-1 p
10. --- 房右邊冇咁多嘢要佢做，因為張\VUgras{長椅} \textsuperscript{(47)} 同佢個
\VUgras{墊} \textsuperscript{(48)} 都擺得啱位，而且因為唔凍，\VUgras{壁爐}
\textsuperscript{(49)} 嘅用具冇乜郁過。把\VUgras{鏟} \textsuperscript{(50)}、對
\VUgras{火鉗} \textsuperscript{(51)}、對\VUgras{柴架} \textsuperscript{(55)} 同塊
\VUgras{擋灰板} \textsuperscript{(56)} 都乾淨；塊\VUgras{擋火屏} \textsuperscript{(54)}
冇畀人郁過，個\VUgras{柴箱} \textsuperscript{(52)} 亦裝滿咗\VUgras{柴}
\textsuperscript{(53)}。

%%K t06-noto-1 noto
\VUnotes{143.2mm}{%
(*) Babo = 細路仔；英文 \textit{baby}，法文 \textit{bébé}，意大利文
\textit{bambino}。%
}

%%K t06-noto-2 noto
\VUnotes{143.2mm}{%
(*) Lambrequino = 法文 \textit{lambrequin}，西班牙文 \textit{lambrequines}，
葡萄牙文 \textit{lambrequins}，連德文同英文都係 \textit{lambrequins}；即係德文、
英文、法文、葡萄牙文同西班牙文都一樣。%
}

%%K t06-c1-11-1 p
11. --- 不過佢仲要掃走個\VUgras{擺鐘} \textsuperscript{(57)}、啲\VUgras{燭台}
\textsuperscript{(58)} 同啲\VUgras{零碎碟} \textsuperscript{(59)} 嘅塵，跟住抹面
\VUgras{鏡} \textsuperscript{(60)}。

%%K t06-c1-12-1 p
12. --- 至於個嬰兒（即係 babo（*））嘅\VUgras{搖籃} \textsuperscript{(61)}，就已經
準備好喇。

%%K t06-tit-3 sub
\VUsaut{5.0mm}
\VUpk{10.2pt}{40}{客廳。}
\VUsaut{3.0mm}

%%K t06-c2-01-1 p
1. --- 而家係客廳。呢間\VUgras{房}好靚。喺右角嘅壁爐，擺住一個精緻嘅\VUgras{小像}
\textsuperscript{(1)} 同兩個靚\VUgras{花瓶} \textsuperscript{(3)}，仲披住一幅絲絨
\VUgras{垂簾} \textsuperscript{(2)}（*）。

%%K t06-c2-02-1 p
2. --- 為咗唔畀爐裏面嘅火星燒著張地氈，人哋喺爐前面擺咗一塊銅嘅\VUgras{擋火星板}
\textsuperscript{(4)}。

%%K t06-c2-03-1 p
3. --- 就喺隔籬，一張雅致嘅女士\VUgras{寫字檯} \textsuperscript{(5)} 開住。
\VUgras{女主人} \textsuperscript{(23)} 就係喺嗰度寫信。

%%K t06-c2-04-1 p
4. --- 隻窗掛住\VUgras{窗紗} \textsuperscript{(6)}，用\VUgras{窗簾帶} \textsuperscript{(8)}
束起，束帶就掛喺\VUgras{掛鈎} \textsuperscript{(7)} 度；上面仲有大幅布簾，頂住一幅
\VUgras{簷飾} \textsuperscript{(9)}。

%%K t06-c2-05-1 p
5. --- 窗嘅凹位裏面有個\VUgras{花盆套} \textsuperscript{(11)}，裝住一棵綠色
\VUgras{植物} \textsuperscript{(10)}，係一棵好靚嘅棕櫚，啲葉伸到差唔多掂到盞
\VUgras{火水燈} \textsuperscript{(12)}。

%%K t06-c2-06-1 p
6. --- 盞燈嘅\VUgras{燈罩} \textsuperscript{(13)} 係瑪利亞整嘅，佢係個迷人嘅
\VUgras{少女} \textsuperscript{(14)}，而家坐喺鋼琴 \textsuperscript{(15)} 前面。佢好直咁坐
喺張\VUgras{琴凳} \textsuperscript{(16)} 上面，練緊一首曲。佢好叻又好細心，睇下佢個
執到井井有條嘅\VUgras{樂譜櫃} \textsuperscript{(17)} 就知。

%%K t06-tit-4 sub
\VUsaut{5.0mm}
\VUpk{10.2pt}{40}{曳仔。}
\VUsaut{3.0mm}

%%K t06-c2-07-1 p
7. --- 佢細佬保羅搵緊一本\VUgras{書} \textsuperscript{(19)}，就喺個\VUgras{書櫃}
\textsuperscript{(18)} 度。佢係個坐唔定又難搞嘅\VUgras{男仔} \textsuperscript{(20)}。
佢攀住個書櫃想攞太高嗰本書，隨時會整跌上面嗰個\VUgras{半身像} \textsuperscript{(21)}。

%%K t06-c2-08-1 p
8. --- 佢趁阿媽同一位朋友傾偈嗰陣做呢件蠢事——嗰位\VUgras{太太} \textsuperscript{(22)}
嚟咗佢屋企住幾日。阿媽坐喺張\VUgras{梳化} \textsuperscript{(24)} 上面，就喺張
\VUgras{大椅} \textsuperscript{(25)} 隔籬，佢個朋友就坐喺張\VUgras{凳}
\textsuperscript{(27)} 上面，我哋淨係見到個\VUgras{椅背} \textsuperscript{(26)}。

%%K t06-c2-09-1 p
9. --- 掛喺牆上面嗰個大\VUgras{相架} \textsuperscript{(30)} 入面係一位親戚嘅
\VUgras{相} \textsuperscript{(29)}。擺喺檯上面嗰本\VUgras{相簿} \textsuperscript{(32)}
入面就收埋咗屋企其他人嘅\VUgras{相片} \textsuperscript{(33)}。啲細路啱啱揭過佢，
因為啲\VUgras{椅} \textsuperscript{(31)} 仲圍住張檯。

%%K t06-c2-10-1 p
10. --- 個瓷\VUgras{中國花瓶} \textsuperscript{(34)} 擺喺把\VUgras{裁紙刀}
\textsuperscript{(35)} 隔籬，佢係人哋喺瑪利亞生日送畀佢嘅，而裝飾住房角嘅嗰塊\VUgras{屏風}
\textsuperscript{(36)} 就係佢阿叔由中國帶返嚟嘅。

%%K t06-c2-11-1 p
11. --- 呢間客廳畀啲靚\VUgras{畫} \textsuperscript{(37)} 令佢生色——嗰啲畫掛喺
\VUgras{間板} \textsuperscript{(42)} 上面——又畀盞\VUgras{吊燈} \textsuperscript{(38)}
照住，盞燈嘅\VUgras{電燈膽} \textsuperscript{(39)} 夜晚開心咁著住，所以睇落好舒服。鋪喺木地板上面
軟綿綿嘅\VUgras{地氈} \textsuperscript{(40)}，同埋咁方便嘅\VUgras{電鐘}
\textsuperscript{(41)}，令佢更加舒適。

%%K t06-tit-5 sub
\VUsaut{3.6mm}
\VUpk{10.2pt}{40}{飯廳。}
\VUsaut{3.0mm}

%%K t06-c3-01-1 p
1. --- 開飯嘅鐘響咗。除咗瑪利亞，成家人都圍埋張檯。飯廳畀一個好好嘅陶
\VUgras{火爐} \textsuperscript{(1)} 焗得好舒服，個爐有條幼幼嘅\VUgras{煙通}
\textsuperscript{(2)}。

%%K t06-c3-02-1 p
2. --- 呢間房冇乜傢俬。沿住幅牆，喺張\VUgras{矮凳} \textsuperscript{(4)} 上面掛住
一個裝飾用嘅細\VUgras{水缸} \textsuperscript{(3)}。就喺隔籬係個\VUgras{碗櫃}
\textsuperscript{(5)}，人哋喺上面擺凍嘅餸、甜品同啲一時未用得著嘅餐具。

%%K t06-c3-03-1 p
3. --- 咁樣，喺上層我哋見到一隻\VUgras{燒雞} \textsuperscript{(7)}、一條\VUgras{魚}
\textsuperscript{(8)} 同一個大\VUgras{杯} \textsuperscript{(10)}，入面裝住
\VUgras{生果} \textsuperscript{(9)}。下面就係\VUgras{芝士} \textsuperscript{(11)}、\VUgras{麵包}
\textsuperscript{(12)} 同個\VUgras{油醋架} \textsuperscript{(13)}，入面有拌沙律要用
嘅嘢：油、醋、\VUgras{鹽} \textsuperscript{(14)} 同\VUgras{胡椒} \textsuperscript{(15)}。
最後，喺最低嗰層有個\VUgras{蛋糕} \textsuperscript{(17)}，就喺個\VUgras{芥辣罌}
\textsuperscript{(16)} 隔籬。

%%K t06-c3-04-1 p
4. --- 飯廳嗰個鐘，人哋叫佢做\VUgras{掛鐘} \textsuperscript{(20)}，樣好特別。佢嵌喺
一個木框裏面，個框仲裝住\VUgras{氣壓計} \textsuperscript{(19)} 同\VUgras{溫度計}
\textsuperscript{(18)}。

%%K t06-tit-6 sub
\VUsaut{4.6mm}
\VUpk{10.2pt}{40}{開晚飯。}
\VUsaut{3.0mm}

%%K t06-c3-05-1 p
5. --- 房入面每幅牆下面都鋪住一行打咗蠟嘅橡木\VUgras{牆板} \textsuperscript{(21)}。
一幅布\VUgras{門簾} \textsuperscript{(22)} 差唔多遮實咗通往騎樓嘅門。食完晚飯，成家
人就會去嗰度飲咖啡。啲\VUgras{杯} \textsuperscript{(24)} 同\VUgras{碟}
\textsuperscript{(25)} 已經擺好喺\VUgras{餐具櫃} \textsuperscript{(23)} 上面。

%%K t06-c3-06-1 p
6. --- 檯上面，天花度裝住一盞\VUgras{吊燈} \textsuperscript{(26)}，佢好光嘅燈夜晚
照住成間飯廳。

%%K t06-c3-07-1 p
7. --- 而家我哋睇下\VUgras{飯檯} \textsuperscript{(27)} 本身。成家人入嚟嗰陣，檯
已經擺好。喺張\VUgras{檯布} \textsuperscript{(28)} 上面，每位食客面前擺住一隻
\VUgras{碟} \textsuperscript{(33)}、一條\VUgras{餐巾} \textsuperscript{(29)}、一隻
\VUgras{匙羹} \textsuperscript{(34)}、一把\VUgras{叉} \textsuperscript{(35)}、一把
\VUgras{刀} \textsuperscript{(36)} 同一隻\VUgras{杯} \textsuperscript{(32)}。仲有
裝酒嘅\VUgras{樽} \textsuperscript{(30)} 同裝水嘅\VUgras{水樽} \textsuperscript{(31)}。

%%K t06-c3-08-1 p
8. --- \VUgras{男工} \textsuperscript{(39)} 首先端上個\VUgras{湯窩} \textsuperscript{(37)}，
入面啲湯仲喺度冒煙。而家佢上緊第一\VUgras{道} \textsuperscript{(38)} 餸，係肉。
跟住佢會上我哋已經講過嗰啲；最後，食完\VUgras{甜品}之後，佢會趁主人過咗去騎樓
嗰陣收埋啲餐具，執返好間飯廳。

%%K t06-c3-08-2 p
\textit{註。} --- \textit{有關食嘢嘅常用詞，喺呢度只係好簡略咁講。喺「酒店」
（第 13 號）同「街市」（第 15 號）兩張圖表度會補足呢張表。}

%%K t06-tit-7 sub
\VUsaut{4.6mm}
\VUpk{10.2pt}{40}{廚房。}
\VUsaut{3.0mm}

%%K t06-c4-01-1 p
1. --- 由半開嘅門望入廚房，我哋見到一片好有畫意嘅凌亂。喺個\VUgras{灶}
\textsuperscript{(1)} 前面——爐入面\VUgras{火} \textsuperscript{(3)} 噼啪響——裝住
一個\VUgras{轉肉架} \textsuperscript{(5)}。一嚿靚\VUgras{燒肉} \textsuperscript{(7)}
\VUgras{串住} \textsuperscript{(6)}，就嚟燒熟。

%%K t06-c4-02-1 p
2. --- 啱啱用完嘅\VUgras{風箱} \textsuperscript{(8)} 同把\VUgras{炭鏟} \textsuperscript{(9)}
仲擺喺地下，啲\VUgras{柴} \textsuperscript{(10)} 都係。

%%K t06-c4-03-1 p
3. --- 爐面執得好整齊：盞\VUgras{燈籠} \textsuperscript{(11)}、個\VUgras{火柴筒}
\textsuperscript{(13)}——入面裝住啲\VUgras{火柴} \textsuperscript{(12)}——個
\VUgras{燭台} \textsuperscript{(14)} 同個\VUgras{矮燭盤} \textsuperscript{(15)}，
上面插住支\VUgras{蠟燭} \textsuperscript{(16)}，全部都擺返好位。不過個大
\VUgras{煤桶} \textsuperscript{(18)} 裝滿\VUgras{煤} \textsuperscript{(17)}——
\VUgras{廚娘} \textsuperscript{(23)} 啱啱落地牢裝返嚟——就仲擺喺廚房中間。

%%K t06-c4-04-1 p
4. --- 的確，呢個可憐嘅女人要趕住整好晏晝飯。而家佢喺個\VUgras{爐頭}
\textsuperscript{(19)} 旁邊，撈緊一味唔可以燒燶咗嘅\VUgras{汁} \textsuperscript{(24)}。

%%K t06-c4-05-1 p
5. --- \VUgras{焗爐} \textsuperscript{(20)} 入面焗緊一個佢為啲細路整嘅餅。爐頭上面
掛住個\VUgras{湯殼} \textsuperscript{(21)} 同個\VUgras{漏殼} \textsuperscript{(22)}。

%%K t06-tit-8 sub
\VUsaut{4.0mm}
\VUpk{10.2pt}{40}{鑊同碗碟。}
\VUsaut{3.0mm}

%%K t06-c4-06-1 p
6. --- 掛喺房尾嘅\VUgras{廚具} \textsuperscript{(25)} 好乾淨。啲靚銅\VUgras{鑊}
\textsuperscript{(26)}、個\VUgras{奶壺} \textsuperscript{(27)}、個\VUgras{隔篩}
\textsuperscript{(28)} 同個\VUgras{刨} \textsuperscript{(29)} 都抹得好乾淨。

%%K t06-c4-07-1 p
7. --- 個\VUgras{鹽罐} \textsuperscript{(30)} 擺喺細櫃上面，就喺個\VUgras{茶壺}
\textsuperscript{(31)} 同裝油嘅\VUgras{大瓶} \textsuperscript{(32)} 隔籬。個
\VUgras{抽屜} \textsuperscript{(33)} 應該係裝住餐具同廚刀。

%%K t06-c4-08-1 p
8. --- 把\VUgras{掃把} \textsuperscript{(34)} 同支\VUgras{雞毛掃} \textsuperscript{(35)}
擺喺一角。廚娘啱啱用過個\VUgras{砧板} \textsuperscript{(36)}；佢將佢擺喺窗邊。

%%K t06-c4-09-1 p
9. --- 一條\VUgras{抹手巾} \textsuperscript{(37)} 掛喺牆上面，就喺個\VUgras{漏斗}
\textsuperscript{(38)} 隔籬。就係喺廚房呢一角，人哋放埋個\VUgras{燒烤架}
\textsuperscript{(39)}、把\VUgras{剁刀} \textsuperscript{(40)} 同塊\VUgras{剁板}
\textsuperscript{(41)}。跟住，喺剁刀上面一塊\VUgras{木板} \textsuperscript{(42)} 上，
擺住幾個\VUgras{罌} \textsuperscript{(43)}、個\VUgras{湯煲} \textsuperscript{(44)}
連佢個\VUgras{蓋} \textsuperscript{(45)}，同個\VUgras{大鑊} \textsuperscript{(46)}。個
\VUgras{煎鑊} \textsuperscript{(47)} 就掛喺附近。

%%K t06-c4-10-1 p
10. --- 呢個角落淨係得個銅\VUgras{水喉} \textsuperscript{(48)} 有啲光。個
\VUgras{洗碗盆} \textsuperscript{(49)}——上面擺住個大\VUgras{水埕} \textsuperscript{(50)}
——同佢個細櫃應該好方便，個櫃入面收住個\VUgras{醋桶} \textsuperscript{(51)}
連佢個\VUgras{喉} \textsuperscript{(52)}、個\VUgras{垃圾桶} \textsuperscript{(53)}
同啲\VUgras{污糟碗碟} \textsuperscript{(54)}。

%%K t06-tit-9 sub
\VUsaut{4.0mm}
\VUpk{10.2pt}{40}{廚房入面第啲嘢。}
\VUsaut{3.0mm}

%%K t06-c4-11-1 p
11. --- 嗰個大\VUgras{盆} \textsuperscript{(55)} 係用嚟洗\VUgras{菜}
\textsuperscript{(58)} 嘅，仲有嗰個生鐵\VUgras{鼎} \textsuperscript{(56)} 擺喺
\VUgras{階磚地} \textsuperscript{(57)} 上面——呢兩樣應該都係收喺嗰個櫃度。

%%K t06-c4-12-1 p
12. --- 廚娘實唔會唔夠火，因為喺檯角度仲有個\VUgras{煤氣爐} \textsuperscript{(59)}，
上面用個細鑊煲緊水。走出嚟嘅\VUgras{水蒸氣} \textsuperscript{(60)} 話畀我哋知啲水
應該滾緊。呢啲水多數係整咖啡用，因為喺檯另一頭就擺住個\VUgras{咖啡壺}
\textsuperscript{(64)} 同個\VUgras{咖啡磨} \textsuperscript{(63)}。

%%K t06-c4-13-1 p
13. --- 個煤氣爐駁住個\VUgras{煤氣掣} \textsuperscript{(62)}，中間係一條長長嘅
\VUgras{膠喉} \textsuperscript{(61)}。

%%K t06-c4-14-1 p
14. --- 嚟幫廚娘手嘅\VUgras{散工女工} \textsuperscript{(66)} 喺度整緊晚飯要用嘅菜。
啲菜擇好洗好之後，佢就會放入個\VUgras{搪瓷盆} \textsuperscript{(65)} 度，等到時候
先至煮。

%%K t06-tit-10 sub
\VUsaut{4.0mm}
{\centering\fontsize{10.2pt}{10.2pt}\selectfont\textls[40]{\scshape 浴室。} \textit{（沖涼房。）}\par}
\VUsaut{3.0mm}

%%K t06-c5-01-1 p
1. --- 我哋淨係仲要好奇多一次，就識晒呢個單位嘅主要房間喇：睇下浴室點裝修。呢樣
唔會耐，因為間房唔大，我哋好快就知個\VUgras{浴缸} \textsuperscript{(1)} 有個好嘅
\VUgras{熱水爐} \textsuperscript{(2)}，個\VUgras{梘盤} \textsuperscript{(3)} 就喺沖涼
嗰個人伸手可及嘅地方，而個\VUgras{毛巾架} \textsuperscript{(5)} 實會令啲\VUgras{毛巾}
\textsuperscript{(4)} 好易乾。

%%K t06-c5-02-1 p
2. --- 呢間房嘅牆鋪住\VUgras{瓷磚} \textsuperscript{(6)}，所以可以裝一個
\VUgras{花灑} \textsuperscript{(7)}，唔使驚用嗰陣濺出嚟嘅水整壞幅牆。仲有個浸腳用
嘅細鋅\VUgras{盆} \textsuperscript{(8)}，令啲設備更加齊全。
\VUsaut{6.0mm}
\VUfilet{22mm}
